\section{Related Literature}
\label{sec:literature}

The paper sits at the intersection of standards policy, private compatibility, and coalition formation. Each component has substantial antecedents. The contribution we claim is correspondingly narrow: the model links a post-formation private adaptation decision to the continuation rents entering governments' comparison of regional and international standardization, and it derives conditions under which that feedback changes sign. This positioning is important because neither one-way compatibility, fixed-cost adoption, standards coalitions, nor regulatory harmonization is new by itself.

\subsection{Private compatibility and costly adaptation}
\label{subsec:lit-private-compatibility}

Private technological responses to incompatibility are a long-standing theme in network economics. \citet{farrellSaloner1992} study converters that connect otherwise incompatible technologies, and \citet{doganogluWright2006} analyze multihoming as a substitute for compatibility. The broader compatibility literature, including \citet{katzShapiro1985} and \citet{besenFarrell1994}, establishes that firms' incentives to interoperate depend on the competitive consequences of a larger compatible network.

More recent work makes the compatibility choice itself strategic. \citet{buccellaEtAl2023} analyze endogenous product compatibility under Cournot competition with network effects and quasi-fixed compatibility costs, including asymmetric compatibility outcomes. In platform settings, \citet{bourreauEtAl2026} study endogenous interoperability with consumer multihoming, while \citet{mottaPeitz2025} analyze strategic denial of interoperability in a dynamic environment. These papers make clear that private compatibility and one-way compatibility are established economic objects. We therefore use outsider-only adoption as a continuation state generated by the model, not as a standalone novelty claim.

The fixed-cost scope logic is also not new. \citet{gorman1985} gives general conditions for economies of scope in the presence of fixed costs, while \citet{choMcCardle2009} show how economies or diseconomies of scope in adoption costs create dependence across multiple technology-adoption decisions. \citet{manentiSomma2008} explicitly study one-way and two-way compatibility in a network-industry entry setting. These parent classes are why the paper treats the outsider-only continuation as supporting structure rather than as a standalone theorem contribution.

Costly multi-standard production and standard adoption likewise have close antecedents. \citet{fischerSerra2000} allow a foreign producer to incur a fixed setup cost when serving markets with different standards. \citet{schmidtSteingress2022} model firms' endogenous adoption of harmonized standards in the presence of sunk adoption costs. In the present model, the fixed cost is standard specific: one investment can cover every market governed by the adopted standard. This creates a scope asymmetry between adopting a two-country bloc standard and adopting a one-country outsider standard. Such fixed-cost scope logic is supporting structure rather than the paper's main result.

The distinction that matters for our analysis is between the existence of a private workaround and its effect on a government's institutional payoff. A converter or additional-standard investment may reduce technological incompatibility while either increasing or decreasing the political return to formal harmonization. Section~\ref{sec:selective-erosion} isolates that second margin through the residual-rent function $R(d,v)$.

\subsection{Standards policy and international harmonization}
\label{subsec:lit-government-standards}

A separate literature studies standards as strategic policy instruments. \citet{gandalShy2001} analyze standardization policy, conversion costs, network effects, and international trade in a three-country environment and allow countries to form standardization unions. \citet{barrettYang2001} show that redesign costs and network effects can generate rational incompatibility with an international product standard, while \citet{klimenko2009} and \citet{klimenkoInteroperability2009} study government policies and agreements over technical compatibility and interoperability. These contributions establish that compatibility affects both firm behavior and government incentives.

Trade-agreement models study related institutional margins. \citet{bagwellStaiger2001} analyze international disciplines over domestic policies, \citet{costinot2008} compares national treatment and mutual recognition for product standards, and \citet{geng2019} studies the same institutional comparison with consumption externalities. \citet{grossmanMcCalmanStaiger2021} examine regulatory convergence when firms tailor product versions to local markets and tailoring generates fixed costs. These papers focus primarily on the design or efficiency of the regulatory arrangement itself. Our formal menu is deliberately simpler; the additional object is a private technological response that occurs after the formal partition has been fixed.

The closest regional-versus-multilateral antecedents are \citet{takaradaEtAl2020} and \citet{kawabataTakarada2021}. \citet{takaradaEtAl2020} compare national, regional, and multilateral standards in a three-country Cournot model using coalition stability and also consider fixed costs of producing under multiple standards. \citet{kawabataTakarada2021} ask whether regional harmonization facilitates or impedes multilateral harmonization. The present paper therefore does not claim novelty for the standards-union menu, the three-country structure, or the regionalism question. The difference is in timing and continuation payoffs: governments choose a formal partition, firms then decide whether to support additional standards, and governments evaluate the resulting product-market outcome without changing formal membership.

Recent work on regulatory diversity reinforces the need for this narrow positioning. \citet{parentiVannoorenberghe2025} show how costly tailoring to country-specific standards can support regulatory blocs. \citet{maggiMrazova2024} study fixed costs of regulatory diversity and the extent to which harmonization can arise without an international regulatory agreement. Together with \citet{schmidtSteingress2022}, these papers show that firm-side adaptation costs and public harmonization are already closely connected. Our contribution is not the broad proposition that private adjustment can substitute for formal cooperation. It is the model-specific sign decomposition identifying when a private response erodes enough of a regional member's rent to strengthen or reverse its relative incentive for broader standardization.

\subsection{Coalition formation and institutional stability}
\label{subsec:lit-coalitions}

Coalition formation provides the institutional application. \citet{hartKurz1983} is a canonical reference for endogenous coalition structures. In the standards literature, \citet{takaradaEtAl2020} explicitly use coalition stability to compare standards regimes. There is also work in which firms themselves form standards coalitions: \citet{economidesSkrzypacz2003} study technical-standards coalition formation in a network industry before oligopoly competition. Outside standards trade, \citet{guoLiuNault2024} study coalition formation when interoperability affects the value of membership.

We do not propose a new coalition solution concept. The canonical application uses exclusive partitions and strict blocking. The role of the product-market model is to generate continuation payoffs for that partition game. R6 robustness analysis also shows why the distinction matters: the symmetric high-fixed-cost stable set depends on strict blocking and symmetry-induced indifference, whereas the intermediate unique-international result is supported by strict payoff differences and survives weak/Pareto blocking in the certified perturbations. The paper therefore separates the government-incentive mechanism from the institution-specific stable-set correspondence.

\subsection{What remains after the parent literatures}
\label{subsec:lit-positioning}

The surrounding ingredients are heavily occupied. Private converters and interoperability are established; fixed-cost adoption and multistandard production are established; governments' standards policies and regional standards unions are established; and firms' costs of regulatory diversity already matter for harmonization. The paper consequently makes no priority claim over those components.

What remains is a conditional interaction among them. In the canonical regional union, the member government's pre-adoption difference relative to international standardization can be written as $\mathscr E-\mathscr D$. After outsider adaptation with residual marginal cost $d$, the same comparison becomes $-\mathscr D+R(d,v)$. The equilibrium-derived term $R(d,v)$ yields distinct success, intermediate, and failure regions rather than a one-directional comparative statement. The incomplete-adaptation open set shows that full cost elimination is not required, while the zero-network, singleton-network, and differentiated-demand exercises identify substantive boundaries to portability.

This combination is the paper's claimed contribution: an explicit mapping from private post-formation adaptation to government continuation incentives, together with analytical conditions under which the mapping reverses sign and concrete models in which the political effect disappears. The coalition-stability results then show one institutional consequence of those continuation-payoff changes rather than serving as the source of the generality claim.
